\providecommand{\BKM}{\mathrm{BKM}}
\providecommand{\ch}{\mathrm{ch}}
\providecommand{\cat}{\mathrm{cat}}
\providecommand{\fib}{\mathrm{fiber}}
\providecommand{\cO}{\mathcal{O}}
\providecommand{\K}{\mathrm{K3}}
\providecommand{\Fix}{\mathrm{Fix}}
\providecommand{\symp}{\mathrm{symp}}
\providecommand{\fr}{\mathrm{fr}}
\providecommand{\CHL}{\mathrm{CHL}}
\providecommand{\GC}{\mathrm{GC}}
\providecommand{\dd}{\mathrm{dd}}
\providecommand{\wt}{\operatorname{wt}}

\section*{Mukai frame refinement of $\kappa_{\BKM}(\Phi_N)$}

\subsection*{Structure-sheaf trace}

A symplectic automorphism $g_N$ of a K3
surface preserves the holomorphic $(2,0)$-form $\Omega_{\K}$, hence acts
trivially on $H^0 \oplus H^2$ of the structure sheaf, giving
$\chi_{g_N}(\cO_{\K}) = 2 = \chi(\cO_{\K})$ for every
$N$. The structure-sheaf trace contains no CHL order dependence. A
calibration through $\chi_{g_N}(\cO_{\K})$ gives the constant row
$c_N(0)=10$ and loses the primitive denominator weights
$5,4,3,2,1$.

The variable datum is the Mukai frame shape of the symplectic action on
the Mathieu $24$-point representation.  If
\[
  \pi(g_N)=\prod_a a^{m_a},
\]
define
\[
  C_{\fr}(g_N):=\sum_a m_a,\qquad
  L_{\top}(g_N):=m_1=\chi(\Fix(g_N)).
\]
Here $L_{\top}$ is the topological Lefschetz number on
$H^*(\K,\mathbb{Q})$, while $C_{\fr}$ is the frame-cycle count entering
the CHL theta component.

\subsection*{CHL denominator row}

\begin{theorem}[Mukai frame refinement of $\kappa_{\BKM}$]
Let $N \in \{1,2,3,4,6\}$, $g_N \in \mathrm{Aut}_{\symp}(\K)$ a Nikulin
symplectic automorphism of order $N$, and let $\psi_N^{\CHL}$ be the
primitive CHL Borcherds input in the Gritsenko and Nikulin
normalisation. Write the index-$1$ theta decomposition
\[
  \psi_N^{\CHL}(\tau, z)
    \;=\; h^{(N)}_0(\tau)\,\theta_{1,0}(\tau, z)
        + h^{(N)}_1(\tau)\,\theta_{1,1}(\tau, z)
\]
(Eichler and Zagier 1985, \S5), and set
\[
  A_N := [q^0]\,h^{(N)}_1(\tau).
\]
Then
\begin{equation}
  c_N^{\CHL}(0) \;=\; C_{\fr}(g_N) \;-\; 2 A_N,
  \qquad
  \kappa_{\BKM}(\Phi_N)
    \;=\; \tfrac12\,c_N^{\CHL}(0).
\end{equation}
The values are
\[
\begin{array}{c c c c c c c}
N & \pi(g_N) & C_{\fr}(g_N) & L_{\top}(g_N) &
A_N & c_N^{\CHL}(0) & \kappa_{\BKM}(\Phi_N)\\[2pt]
1 & 1^{24} & 24 & 24 & 7 & 10 & 5\\
2 & 1^8 2^8 & 16 & 8 & 4 & 8 & 4\\
3 & 1^6 3^6 & 12 & 6 & 3 & 6 & 3\\
4 & 1^4 2^2 4^4 & 10 & 4 & 3 & 4 & 2\\
6 & 1^2 2^2 3^2 6^2 & 8 & 2 & 3 & 2 & 1
\end{array}
\]
\end{theorem}

\begin{proof}
Mukai 1988, Theorem~0.6, and Nikulin 1979, \S4, embed finite
symplectic K3 automorphism groups in the Mathieu permutation
representation.  The five CHL orders have frame shapes
\[
  1^{24},\quad 1^8 2^8,\quad 1^6 3^6,\quad
  1^4 2^2 4^4,\quad 1^2 2^2 3^2 6^2 .
\]
Each frame shape has total degree $24$.  The topological fixed-point
count is the coefficient $m_1$, giving
\[
  L_{\top}(g_N)=(24,8,6,4,2).
\]
The theta refinement uses the cycle count
\[
  C_{\fr}(g_N)=\sum_a m_a=(24,16,12,10,8).
\]

The $q^0$ part of the CHL input has the form
\[
  \psi_N^{\CHL}\big|_{q^0}
    \;=\; A_N\zeta^{-1}+c_N^{\CHL}(0)+A_N\zeta .
\]
Evaluation at $\zeta=1$ gives
$C_{\fr}(g_N)=c_N^{\CHL}(0)+2A_N$ in the CHL frame normalisation.
The theta corrections
\[
  A_N=(7,4,3,3,3)
\]
are the leading odd-component coefficients in the Gritsenko and
Nikulin paramodular input (Gritsenko 1999, \S6; Eguchi and Hikami
2011, Table~1).  Substitution gives
\[
  c_N^{\CHL}(0)=(10,8,6,4,2).
\]
Borcherds' singular theta lift identifies the modular weight with half
the discriminant-zero constant coefficient (Borcherds 1998,
Theorem~13.3):
\[
  \wt(\Phi_N)=c_N^{\CHL}(0)/2.
\]
This proves
\[
  \kappa_{\BKM}(\Phi_N)=(5,4,3,2,1)
\]
on the primitive CHL denominator branch.
\end{proof}

\subsection*{Non-CHL Nikulin boundary}

The order-indexed Nikulin extension has three further automorphic
rows.  The same frame calculation gives
\[
\begin{array}{c c c c c c c}
N & \pi(g_N) & C_{\fr}(g_N) & L_{\top}(g_N) &
A_N & c_N(0) & \wt(\Phi_N)\\[2pt]
5 & 1^4 5^4 & 8 & 4 & 2 & 4 & 2\\
7 & 1^3 7^3 & 6 & 3 & 2 & 2 & 1\\
8 & 1^2 2^1 4^1 8^2 & 6 & 2 & 2 & 2 & 1
\end{array}
\]
These are Borcherds automorphic weights after the required cover,
input, and normalisation choices.  The notation
$\kappa_{\BKM}(\Phi_N)$ is used for the primitive CHL denominator row;
the boundary rows require an additional denominator algebra and
Calabi-Yau host before carrying the same geometric meaning.

\subsection*{Diagonal-divisor catalogue}

The Gritsenko and Clery diagonal-divisor catalogue is indexed by
paramodular triples.  The cyclic symplectic order is a different index.
Its constants are
\[
\begin{array}{c c c c c c c c c}
\mathrm{row} & 1&2&3&4&5&6&7&8\\[2pt]
\wt & 5&2&3&1&2&1/2&3/2&1\\
c_{\dd}^{\GC}(0) & 10&4&6&2&4&1&3&2
\end{array}
\]
with $\wt=c_{\dd}^{\GC}(0)/2$ row by row.  This table is a separate
paramodular atlas.  Agreement with the CHL branch at the $\Delta_5$
entry fixes the shared level-one normalisation.

\subsection*{References}

Mukai 1988, \emph{Invent.\ Math.} 94, 183 to 221, Table~2.
Atiyah and Bott 1968, \emph{Ann.\ of Math.} 88, 451 to 491,
Theorem~4.12.
Nikulin 1979, \emph{Trans.\ Moscow Math.\ Soc.} 38, \S4.
Eichler and Zagier 1985, \emph{Theory of Jacobi Forms}, Progress in
Math.\ 55, \S5.
Borcherds 1998, \emph{Invent.\ Math.} 132, 491 to 562, Theorem~13.3.
Eguchi, Ooguri, and Tachikawa 2011, \emph{Exp.\ Math.} 20, Table~1.
Eguchi and Hikami 2011, \emph{Phys.\ Lett.} B 694, Table~1.
Cheng, Duncan, and Harvey 2014, \emph{Commun.\ Num.\ Theory Phys.} 8,
Table~4.
Gritsenko 1999, arXiv:math/9906191.
Gritsenko and Nikulin 1998, \emph{Amer.\ J.\ Math.} 119.
Gritsenko and Clery 2013, \emph{J.\ Reine Angew.\ Math.} 678.
Lorgat 2020, \emph{A Borcherds lift of $\phi_{0,1}$}, \S2.
